% !TeX root = ../main.tex
\section{Покрытие исходника и исправления}
\subsection*{Соответствие страницам рукописи}
Страница 1 --- декоративная обложка; вместо неё использован титульный
заголовок в стиле остальных проектов. Все самостоятельные решения
со страниц 2--42 включены. Порядок задач внутри разделов сохранён.
\begin{center}
\begin{tabular}{r|p{12.5cm}}
Страницы & Задачи\\\hline
2--3 & С. 1:13; С. 2:7, 10\\
4--6 & С. 2:44; № 55 в рукописи; Ф. 62\\
6--9 & С. 2:59, 73, 80; С. 3:25\\
9--11 & С. 3:59, 68\\
12--14 & С. 3:94; Ф. 146; С. 4:4, 20\\
15--16 & С. 4:20, 59; Д. 1; начало С. 7:1\\
16--19 & С. 7:1, 28, 5\\
20--23 & С. 7:5, 46, 65, а\\
24--27 & Ф. 505; Д. 2; Ф. 225, а, г\\
28--30 & Ф. 228, в, г; Ф. 229; Ф. 230\\
31--33 & Ф. 231; Ф. 233; Д. 3, а, б; начало Ф. 278\\
33--37 & Ф. 278, 287, 288\\
38--40 & С. 6:7; Д. 5, а, б\\
41--42 & Д. 6
\end{tabular}
\end{center}
\subsection*{Существенные исправления и дополнения}
\begin{enumerate}
\item В № 55 добавлено полное семейство решений со склейкой
кубических ветвей и нулевого участка. Два ответа рукописи не исчерпывают
всех решений задачи Коши.
\item В С. 2:7 запрет $\cos x=0$ отделён от ограничений на промежуточное
деление: гладкое продолжение через такие точки возможно.
\item В С. 2:59 константа после потенцирования обязана быть положительной.
Уточнён смысл вертикальной кривой $x=0$.
\item В С. 2:73 и 80 пояснён выбор непрерывной ветви угла;
для задачи Коши С. 2:80 указан максимальный интервал графика.
\item В С. 3:25 и 59 описано поведение в вырожденной точке $x=0$.
В С. 3:68 цепочка замен упрощена переходом к функции $x(y)$.
\item В С. 4:20 исправлен знак в уравнении на интегрирующий множитель:
$\mu'/\mu=-2/y$.
\item В С. 7:5 исправлена промежуточная линейная формула для $1/y$.
В С. 7:46 восстановлены все постоянные решения.
\item В С. 7:65, а исправлен знак при $3\dot z$ и проверено продолжение
явного ответа через нуль. Для Д. 2 отмечена неединственность продолжения
через вырожденную точку, несмотря на локальную единственность около $x=1$.
\item В С. 7:28 указан правильный интервал задачи Коши;
в Ф. 505 явно выбрана ветвь ответа по начальным данным.
\item В Ф. 225 исправлена область непрерывности $\tan y$.
В Ф. 228 различены достаточный критерий единственности и необходимые
условия; фазовое пространство уравнения третьего порядка имеет четыре координаты
с учётом независимой переменной.
\item В Ф. 233 достаточность условия $n\ge5$ доказана примером
$y^{(n)}=0$, а не только отсутствием противоречия.
\item В Д. 3, а область решения задачи Коши исправлена с
$\R\setminus\{2\}$ на $(-\infty,2)$.
\item В Ф. 288 исправлен знак константы во второй параметрической
координате и добавлено потерянное решение $y=0$.
\item В Ф. 278, 287 и С. 6:7 описаны допустимые $C^1$-склейки
с особыми решениями.
\item В Д. 5, а исключён ложный второй ответ: до $x=6$ он идёт
по параболе и при $x=5$ даёт $25/4$, а не $6$.
В Д. 5, б единственный ответ рукописи дополнен всеми продолжениями
с произвольной точкой выхода $b\ge4$.
\item В Д. 6 объединены два одинаковых семейства со знаками
$\pm C$, добавлены задержки на $y=1$ и обосновано отсутствие
склеек с $y=0$ за конечное время.
\end{enumerate}
\subsection*{Как использованы приложенные материалы}
Основной источник --- \texttt{дз1\_чат\_кураторов.pdf}, 42 страницы.
Формулы задач с обозначением «С.» сверялись с приложенным сборником
под редакцией В.\,К.\,Романко там, где требовалось уточнить постановку
(в частности, ортогональные траектории С. 2:44 и данные С. 7:65).
В № 55 и Д. 5 сохранены и явно отмечены расхождения нумерации или
ссылок между рукописью и другими приложенными материалами.
Конспект курса и программа не использовались как указание добавлять
в документ другие задачи. Балльно-рейтинговая таблица не относится
к набираемым решениям и в текст не включена.
